% 第1章 1.5.4：三种收敛方式的反例对照。
% (a) x^n（n=2,8,30）：逐点趋于 1_{1}，L^p 范数 (np+1)^{-1/p}->0，上确界恒为 1。
% (b) 窄尖峰 g_n=n 1_{(0,1/n)}（n=2,4,8）：处处趋于 0，||g_n||_1=1。
% (c) 打字机序列：第 k 层依次为 [j2^{-k},(j+1)2^{-k}) 的指示函数（k=0..3，共 15 项），横条为支撑；
%     竖虚线 x=0.3 在每层恰好命中一项。
\begin{tikzpicture}[font=\figurefont\footnotesize]
\pgfplotsset{every axis/.style={width=40mm,height=38mm,axis lines=left,
  x axis line style={draw=FigureStroke,line width=.6pt},y axis line style={draw=FigureStroke,line width=.6pt},
  tick style={draw=FigureStroke},tick label style={font=\figurefont\scriptsize,text=FigureInk},
  label style={font=\figurefont\scriptsize},
  xlabel style={at={(axis description cs:1,0)},anchor=north west,yshift=1mm},
  title style={font=\figurefont\footnotesize},samples=200}}
\begin{axis}[name=A,xmin=0,xmax=1,ymin=0,ymax=1.08,xtick={0,1},ytick={0,1},xlabel={$x$},title={(a) $x^n$}]
\addplot[draw=FigureBlue,line width=1.1pt,domain=0:1] {x^2};
\addplot[draw=FigureGreen,line width=1.1pt,domain=0:1] {x^8};
\addplot[draw=FigureYellow,line width=1.1pt,domain=0:1,samples=300] {x^30};
\node[font=\figurefont\scriptsize,text=FigureBlue] at (axis cs:0.42,0.40) {$2$};
\node[font=\figurefont\scriptsize,text=FigureGreen] at (axis cs:0.70,0.30) {$8$};
\node[font=\figurefont\scriptsize,text=FigureYellow] at (axis cs:0.84,0.12) {$30$};
\end{axis}
\begin{axis}[name=B,at={(A.east)},anchor=west,xshift=9mm,xmin=0,xmax=1,ymin=0,ymax=8.6,xtick={0,0.5,1},xticklabels={$0$,$\tfrac12$,$1$},ytick={2,4,8},xlabel={$x$},title={(b) $n\mathbf 1_{(0,1/n)}$}]
\addplot[draw=FigureBlue,line width=1.1pt] coordinates {(0,2) (0.5,2) (0.5,0) (1,0)};
\addplot[draw=FigureGreen,line width=1.1pt] coordinates {(0,4) (0.25,4) (0.25,0) (1,0)};
\addplot[draw=FigureYellow,line width=1.1pt] coordinates {(0,8) (0.125,8) (0.125,0) (1,0)};
\end{axis}
\begin{axis}[at={(B.east)},anchor=west,xshift=9mm,xmin=0,xmax=1,ymin=0.3,ymax=15.7,y dir=reverse,xtick={0,0.3,1},xticklabels={$0$,$x_0$,$1$},ytick={1,3,7,15},
  xlabel={$x$},title={(c) 打字机序列}]
\addplot[draw=FigureMuted,line width=.6pt,dashed] coordinates {(0.3,0.3) (0.3,15.7)};
\addplot[draw=FigureBlue,line width=1.6pt] coordinates {(0,1) (1,1)};
\addplot[draw=FigureGreen,line width=1.6pt] coordinates {(0,2) (0.5,2)};
\addplot[draw=FigureGreen,line width=1.6pt] coordinates {(0.5,3) (1,3)};
\foreach \j in {0,...,3} {\edef\tmp{\noexpand\addplot[draw=FigureYellow,line width=1.6pt] coordinates {(\j*0.25,\j+4) (\j*0.25+0.25,\j+4)};}\tmp}
\foreach \j in {0,...,7} {\edef\tmp{\noexpand\addplot[draw=FigurePurple,line width=1.6pt] coordinates {(\j*0.125,\j+8) (\j*0.125+0.125,\j+8)};}\tmp}
\end{axis}
\end{tikzpicture}
